% GENERATED FILE. Edit canonical lessons, then run npm run generate:books.
% canonical-source-hash: fnv1a64:fd486876697d3612
% canonical-lessons: LA-C162-nuper, LA-C162-nudiustertius, LA-C162-olim, LA-C162-antea, LA-C162-accido

\chapter{Recently, the Day before Yesterday, Once, Before, To Happen}
\label{ch:la-recently-the-day-before-yesterday-once-before-to-happen}
\hlchaptermodality{162}

\begin{chapteropening}
\textbf{By the end of this chapter:} \emph{I can say what happened, and when: recently, the day before yesterday, long ago.}

Complete the last lesson of chapter 162: I can say what happened, and when: recently, the day before yesterday, long ago.
\end{chapteropening}

\section[nūper]{nūper — recently}

\label{lesson:LA-C162-nuper}



\begin{quote}
Before the new one: say the Latin for which, then the Latin for and not.
\end{quote}

\subsection*{You'll want to know: nūper}
\textbf{nūper} — \textquotedblleft{}recently\textquotedblright{}. Latin's past tense, the perfect, has its own endings: \textbf{vīdī} — \textquotedblleft{}I saw\textquotedblright{}, \textbf{vēnī} — \textquotedblleft{}I came\textquotedblright{}. \textbf{Nūper librum lēgī.} — \textquotedblleft{}I recently read a book.\textquotedblright{}

\begin{grammarlens}[title={What you've built}]
One more word for this chapter.
\end{grammarlens}

\subsection*{Your turn}
\begin{itemize}
\raggedright
  \item \emph{Say it:} \emph{nūper}
  \item \emph{Say it:} \emph{nūper}, once more
  \item \emph{Say it:} say \emph{nūper}
  \item \emph{Recall:} say the Latin for which, then the Latin for and not, then say \emph{nūper} again
\end{itemize}

\begin{tcolorbox}[breakable,colback=teal!4,colframe=teal!35!black,title={Before you move on}]
What does nūper mean? (\textquotedblleft{}Recently\textquotedblright{}.) Say it once more.
\end{tcolorbox}
\section[nudiustertius]{nudiustertius — the day before yesterday}

\label{lesson:LA-C162-nudiustertius}



\begin{quote}
Before the new one: say the Latin for and not, then the Latin for recently.
\end{quote}

\subsection*{You'll want to know: nudiustertius}
\textbf{nudiustertius} — \textquotedblleft{}the day before yesterday\textquotedblright{}. It is \textbf{nunc diēs tertius}, \textquotedblleft{}now (is) the third day\textquotedblright{}: counting today, yesterday is the second. \textbf{Nudiustertius epistulam scrīpsī.} — \textquotedblleft{}The day before yesterday I wrote a letter.\textquotedblright{}

\begin{grammarlens}[title={What you've built}]
One more word for this chapter.
\end{grammarlens}

\subsection*{Your turn}
\begin{itemize}
\raggedright
  \item \emph{Say it:} \emph{nudiustertius}
  \item \emph{Say it:} \emph{nudiustertius}, once more
  \item \emph{Say it:} say \emph{nudiustertius}
  \item \emph{Recall:} say the Latin for and not, then the Latin for recently, then say \emph{nudiustertius} again
\end{itemize}

\begin{tcolorbox}[breakable,colback=teal!4,colframe=teal!35!black,title={Before you move on}]
What does nudiustertius mean? (\textquotedblleft{}The day before yesterday\textquotedblright{}.) Say it once more.
\end{tcolorbox}
\section[ōlim]{ōlim — once, long ago; one day}

\label{lesson:LA-C162-olim}



\begin{quote}
Before the new one: say the Latin for recently, then the Latin for the day before yesterday.
\end{quote}

\subsection*{You'll want to know: ōlim}
\textbf{ōlim} — \textquotedblleft{}once, long ago; one day\textquotedblright{}. It points to a time far off, usually long past: \textbf{Ōlim Rōma parva erat.} — \textquotedblleft{}Long ago Rome was small.\textquotedblright{} \textbf{erat} is \textquotedblleft{}was\textquotedblright{}, from \textbf{sum}.

\begin{grammarlens}[title={What you've built}]
One more word for this chapter.
\end{grammarlens}

\subsection*{Your turn}
\begin{itemize}
\raggedright
  \item \emph{Say it:} \emph{ōlim}
  \item \emph{Say it:} \emph{ōlim}, once more
  \item \emph{Say it:} say \emph{ōlim}
  \item \emph{Recall:} say the Latin for recently, then the Latin for the day before yesterday, then say \emph{ōlim} again
\end{itemize}

\begin{tcolorbox}[breakable,colback=teal!4,colframe=teal!35!black,title={Before you move on}]
What does ōlim mean? (\textquotedblleft{}Once, long ago; one day\textquotedblright{}.) Say it once more.
\end{tcolorbox}
\section[anteā]{anteā — before, previously}

\label{lesson:LA-C162-antea}



\begin{quote}
Before the new one: say the Latin for the day before yesterday, then the Latin for once.
\end{quote}

\subsection*{You'll want to know: anteā}
\textbf{anteā} — \textquotedblleft{}before, previously\textquotedblright{}. It is \textbf{ante}, \textquotedblleft{}before\textquotedblright{}, with \textbf{eā}. \textbf{Anteā Rōmae habitāvī.} — \textquotedblleft{}Before, I lived in Rome.\textquotedblright{} \textbf{habitāvī} is the perfect of \textbf{habitō}.

\begin{grammarlens}[title={What you've built}]
One more word for this chapter.
\end{grammarlens}

\subsection*{Your turn}
\begin{itemize}
\raggedright
  \item \emph{Say it:} \emph{anteā}
  \item \emph{Say it:} \emph{anteā}, once more
  \item \emph{Say it:} say \emph{anteā}
  \item \emph{Recall:} say the Latin for the day before yesterday, then the Latin for once, then say \emph{anteā} again
\end{itemize}

\begin{tcolorbox}[breakable,colback=teal!4,colframe=teal!35!black,title={Before you move on}]
What does anteā mean? (\textquotedblleft{}Before, previously\textquotedblright{}.) Say it once more.
\end{tcolorbox}
\section[accidō, accidere]{accidō, accidere — to happen}

\label{lesson:LA-C162-accido}



\begin{quote}
Before the new one: say the Latin for once, then the Latin for before.
\end{quote}

\subsection*{You'll want to know: accidō, accidere}
\textbf{accidō, accidere} — \textquotedblleft{}to happen\textquotedblright{}. \textbf{Quid accidit?} — \textquotedblleft{}What happened?\textquotedblright{} Its perfect is spelled like its present in this form: \textbf{accidit}. It is \textbf{ad} + \textbf{cadō}, \textquotedblleft{}to fall towards\textquotedblright{}, the source of English \emph{accident}.

\begin{grammarlens}[title={What you've built}]
One more word for this chapter.
\end{grammarlens}

\subsection*{Your turn}
\begin{itemize}
\raggedright
  \item \emph{Say it:} \emph{accidō, accidere}
  \item \emph{Say it:} \emph{accidō, accidere}, once more
  \item \emph{Say it:} say \emph{accidō, accidere}
  \item \emph{Recall:} say the Latin for once, then the Latin for before, then say \emph{accidō, accidere} again
\end{itemize}

\begin{tcolorbox}[breakable,colback=teal!4,colframe=teal!35!black,title={Before you move on}]
What does accidō, accidere mean? (\textquotedblleft{}To happen\textquotedblright{}.) Say it once more.
\end{tcolorbox}
